% Propietario: output/LEY9_SINTESIS_300_SUCESORA_SIN_REGRESIONES_20260904/
% manuscrito/sections/autonomia/03_prolongacion_tpk.tex:354--496.
% Microetapa CG-007. Exposicion autocontenida de dominios y prueba.
\section{Coinductive boundary and stabilization of compatible extensions}
\label{scope:g9-frontera}

The initial emission has image $\mathcal W_6$ of cardinality $243$ within the space of $729$ six-trit words. Its extensions preserve the enriched state: word, position, direction, sheet, regime, quotients, carry, boundary, incidences and chronological register. For $r\ge1$, denote by $\widehat{\mathcal W}_{6r}$ the set of admissible states of length $6r$ actually obtained through TPK transitions, with their compatible prefixes and survival coordinate. The word projection takes values in $(\mathbb F_3^6)^r$.

Truncation $t_{r+1,r}$ retains the enriched prefix of length $6r$. The extension relation is
\[
 \mathcal P_r=\{(x,y):x\in\widehat{\mathcal W}_{6r},\
 y\in\widehat{\mathcal W}_{6(r+1)},\ t_{r+1,r}(y)=x\}.
\]
The relations use admissible states produced by the TPK; equality of prefixes expresses their compatibility. For $s>r$, write $t_{s,r}=t_{r+1,r}\circ\cdots\circ t_{s,s-1}$. Preservation of fields requires $t_{u,r}=t_{s,r}t_{u,s}$ when $u>s>r$.

\begin{proposition}[Restitution of the antecedent by truncation]
If $p_r$ is an extension section whose graph is contained in $\mathcal P_r$, then $t_{r+1,r}p_r=I$. For a sequence of such sections, composing their truncations recovers the full antecedent state.
\end{proposition}
\begin{proof}
Graph inclusion means, for every $x$ in the section's domain, $t_{r+1,r}(p_r(x))=x$. When two sections are composed, the deeper pair cancels first, followed by the preceding pair. Induction proves the formula for any finite number of stages. Equality concerns enriched states, hence preserves all their fields and prefix relations.
\end{proof}

\subsection{Candidates at the first boundary and opening relation}

Let $\widetilde{\mathcal R}_{36}$ be the set of admissible states of depth $36$ with defined boundary register and edge memory. For $x\in\widetilde{\mathcal R}_{36}$, the set $\operatorname{Cand}_6(x)$ consists of states $y$ of depth $42$ obtained by a typed trajectory
\[
 x=x_{36}\xrightarrow{\mathcal U_{36}}x_{37}
 \xrightarrow{\mathcal U_{37}}\cdots
 \xrightarrow{\mathcal U_{41}}x_{42}=y,
 \qquad \tau_{42,36}(y)=x.
\]
Each candidate's visible tail has six trits, so its word projection has at most $729$ values. The cardinality of this projection is distinct from the cardinality of states with complete registers.

The predicate $g_{9,x}(y)$ requires all prefixes to be compatible and admissible continuations to exist at arbitrary depth. The effective boundary and its opening are defined by
\begin{align*}
 \mathcal R_{36}
 &=\{x\in\widetilde{\mathcal R}_{36}:
      \exists y\in\operatorname{Cand}_6(x),\ g_{9,x}(y)=1\},\\
 G_9&=\{(x,y):x\in\mathcal R_{36},\
                y\in\operatorname{Cand}_6(x),\ g_{9,x}(y)=1\}.
\end{align*}
The construction preserves the truncation square $\tau_{42,36}(y)=x$ and gives
\[
 G_9\subseteq\mathcal U_{41}\circ\cdots\circ\mathcal U_{36}
 \subseteq\widehat{\mathcal X}_{36}\times\widehat{\mathcal X}_{42}.
\]
Each element of $\mathcal R_{36}$ has a surviving candidate by definition. The relation $G_9$ opens a six-trit block; the holonomy $\Gamma_9$ composes nine elementary transitions. The first object is a relation between depth levels, and the second is the calendar return with its memory increment.

\subsection{Existence of a complete history}

\begin{theorem}[Stabilization with memory preservation]
Suppose the admissible prefixes arise from a finite nonempty set of initial conditions and form a finitely branching system with states at every depth, whose truncations satisfy the preceding compatibility and preserve the state fields. Suppose also that nonadic transport is natural with respect to truncation and that the return identities $g_{k+9}=g_k$, $m_{k+9}=m_k+1$ hold on admitted trajectories. Then a compatible infinite history exists. Along it, each return preserves phase and increases integer memory by one unit.
\end{theorem}
\begin{proof}
Each prefix is organized by length and linked to its immediate extensions. If there are several initial states, a formal root with that finite set of successors is added. Non-emptiness at every depth provides arbitrarily long extensions from this root.

A vertex admitting arbitrarily long extensions has at least one successor with the same property. Indeed, if each of its finitely many successors had a finite depth bound, the maximum of those bounds, increased by one, would bound all extensions of the vertex. This contradicts the initial property. A successor with unbounded extension is chosen successively; the finite ordering of words and registers fixes the choice when this encoding is available. Iteration produces a chain with a prefix at every level.

The extension relation guarantees that every immediate truncation restores its antecedent. Compatibility under composition proves the same equality between any pair of depths. Naturality allows the return identity to be applied to every prefix where it is defined. Its two components give, respectively, preserved phase and the unit memory increment. After $s$ turns, $g_{k+9s}=g_k$ and $m_{k+9s}=m_k+s$, by induction.
\end{proof}

The proof establishes existence from finite branching, extension to every depth and compatibility. The survival predicate expresses this property of the complete trajectory. Individualizing a concrete regional section additionally uses that region's reader and functional rule, developed in its own chain. These two results retain their respective domains and compose when the selected section satisfies the preceding extension relations.
